\setcounter{section}{10}






  \zwischenueberschrift{Übungsaufgaben}

\inputaufgabegibtloesung
{}
{





Es sei $K$ ein
\definitionsverweis {Körper}{}{}
und sei $A$ eine kommutative
$K$-\definitionsverweis {Algebra}{}{,}
die als
$K$-\definitionsverweis {Modul}{}{}
\definitionsverweis {endlich}{}{}
sei. Zeige, dass ein Element

\mavergleichskette
{\vergleichskette
{ f
}
{ \in }{ A
}
{  }{
}
{    }{
}
{   }{
}
}
{}{}{}
genau dann eine
\definitionsverweis {Einheit}{}{}
ist, wenn es ein
\definitionsverweis {Nichtnullteiler}{}{}
ist.





}
{} {}

\inputaufgabe
{}
{





Es seien $K$ und $L$ Körper, es sei

\mavergleichskette
{\vergleichskette
{ K
}
{ \subseteq }{ L
}
{  }{
}
{    }{
}
{   }{
}
}
{}{}{}
eine
\definitionsverweis {endliche Körpererweiterung}{}{}
und sei $A$,

\mavergleichskette
{\vergleichskette
{ K
}
{ \subseteq }{ A
}
{ \subseteq }{ L
}
{    }{
}
{   }{
}
}
{}{}{,}
ein Zwischenring. Zeige, dass dann $A$ ebenfalls ein Körper ist.





}
{} {}

\inputaufgabe
{}
{





Es sei

\mavergleichskette
{\vergleichskette
{ R
}
{ \subseteq }{ S
}
{  }{
}
{    }{
}
{   }{
}
}
{}{}{}
eine
\definitionsverweis {endliche Ringerweiterung}{}{}
und sei

\mavergleichskette
{\vergleichskette
{ f
}
{ \in }{ R
}
{  }{
}
{    }{
}
{   }{
}
}
{}{}{.}
Zeige: Wenn $f$, aufgefasst in $S$, eine
\definitionsverweis {Einheit}{}{}
ist, dann ist $f$ eine Einheit in $R$.





}
{} {}

\inputaufgabe
{}
{





Es sei $R$ ein
\definitionsverweis {kommutativer Ring}{}{} und $M$ ein
$R$-\definitionsverweis {Modul}{}{.} Dann ist $M$ genau dann
\definitionsverweis {noethersch}{}{,}
wenn jede aufsteigende Kette

\mavergleichskettedisp
{\vergleichskette
{ M_0
}
{ \subseteq} { M_1
}
{ \subseteq} { M_2
}
{ \subseteq} { \ldots
}
{ } {
}
}
{}{}{}
von
$R$-\definitionsverweis {Untermoduln}{}{}
\definitionsverweis {stationär}{}{}
wird.





}
{} {}

\inputaufgabe
{}
{





Es sei

\mavergleichskette
{\vergleichskette
{ f
}
{ \in }{ R
}
{  }{
}
{    }{
}
{   }{
}
}
{}{}{}
ein
\definitionsverweis {Nichtnullteiler}{}{}
in einem
\definitionsverweis {kommutativen Ring}{}{}
$R$. Zeige, dass dies zu einer
\definitionsverweis {kurzen exakten Sequenz}{}{}

\mathdisp {0\stackrel{}{\longrightarrow}R
\stackrel{\cdot f}{\longrightarrow}R

 \mathdisplaybruch\stackrel{}{\longrightarrow}R/(f)
\stackrel{}{\longrightarrow}0} {  }

von
$R$-\definitionsverweis {Moduln}{}{}
führt.





}
{} {}

\inputaufgabegibtloesung
{}
{





Es sei $R$ ein
\definitionsverweis {kommutativer Ring}{}{}
und seien

\mavergleichskette
{\vergleichskette
{ I,J
}
{ \subseteq }{ R
}
{  }{
}
{    }{
}
{   }{
}
}
{}{}{}
\definitionsverweis {Ideale}{}{.}
Zeige, dass die Sequenz

\mathdisp {0\stackrel{}{\longrightarrow}R/ I \cap J
\stackrel{}{\longrightarrow}R/I \times R/J

 \mathdisplaybruch\stackrel{}{\longrightarrow}R/I+J
\stackrel{}{\longrightarrow}0} {  }

mit
\mathl{r \mapsto (r,r)}{} und
\mathl{(s,t) \mapsto s-t}{}
\definitionsverweis {exakt}{}{}
ist.





}
{} {}

\inputaufgabe
{}
{





Es sei $R$ ein
\definitionsverweis {kommutativer Ring}{}{}
und sei $N$ ein
$R$-\definitionsverweis {Modul}{}{}
mit
$R$-\definitionsverweis {Untermoduln}{}{}

\mavergleichskette
{\vergleichskette
{ L
}
{ \subseteq }{ M
}
{ \subseteq }{ N
}
{    }{
}
{   }{
}
}
{}{}{.}
Zeige, dass die
\definitionsverweis {Restklassenmoduln}{}{}
durch die
\definitionsverweis {kurze exakte Sequenz}{}{}

\mathdisp {0\stackrel{}{\longrightarrow}M/L
\stackrel{}{\longrightarrow}N/L

 \mathdisplaybruch\stackrel{}{\longrightarrow}N/M
\stackrel{}{\longrightarrow}0} {  }

miteinander in Beziehung stehen.





}
{} {}

\inputaufgabe
{}
{





Es sei $R$ ein
\definitionsverweis {kommutativer Ring}{}{}
und sei
\maabb {\varphi} { M } { N
} {}
ein
$R$-\definitionsverweis {Modulhomomorphismus}{}{}
zwischen den
$R$-\definitionsverweis {Moduln}{}{}
\mathkor {} {M} {und} {N} {.}
Zeige, dass dies zu einer
\definitionsverweis {kurzen exakten Sequenz}{}{}

\mathdisp {0 \longrightarrow \operatorname{kern} \varphi  \longrightarrow M  \longrightarrow \operatorname{bild} \varphi \longrightarrow  0} {  }

führt.





}
{} {}







Es sei $R$ ein
\definitionsverweis {kommutativer Ring}{}{} und $M$ ein
$R$-\definitionsverweis {Modul}{}{.} Der $R$-Modul

\mavergleichskettedisp
{\vergleichskette
{ {  M   }^{ * }
}
{ =} { \operatorname{Hom}_{  R } \left(   M  ,  R   \right)
}
{ } {
}
{ } {
}
{ } {
}
}
{}{}{}
heißt der
\definitionswort {duale Modul}{}
zu $M$.







\inputaufgabegibtloesung
{}
{





Es sei $R$ ein
\definitionsverweis {kommutativer Ring}{}{} und sei

\mathdisp {0 \, \stackrel{  }{\longrightarrow} \,  L   \, \stackrel{  }{ \longrightarrow}   \,  M   \, \stackrel{  }{\longrightarrow} \,  N \,\stackrel{  } {\longrightarrow}  \, 0} {  }

eine
\definitionsverweis {kurze exakte Sequenz}{}{}
von
$R$-\definitionsverweis {Moduln}{}{}

\mathl{L,M,N}{.} Zeige, dass dies zu einer exakten Sequenz

\mathdisp {0 \longrightarrow {  N   }^{ * } \longrightarrow {  M   }^{ * }  \longrightarrow {  L   }^{ * }} {  }

der
\definitionsverweis {dualen Moduln}{}{}
führt.





}
{} {}

\inputaufgabe
{}
{





Es sei $K$ ein
\definitionsverweis {Körper}{}{} und sei

\mathdisp {0 \, \stackrel{  }{\longrightarrow} \,  L   \, \stackrel{  }{ \longrightarrow}   \,  M   \, \stackrel{  }{\longrightarrow} \,  N \,\stackrel{  } {\longrightarrow}  \, 0} {  }

eine
\definitionsverweis {kurze exakte Sequenz}{}{}
von
$K$-\definitionsverweis {Vektorräumen}{}{}

\mathl{L,M,N}{.} Zeige, dass dies zu einer kurzen exakten Sequenz

\mathdisp {0 \, \stackrel{  }{\longrightarrow} \,  {  N   }^{ * }   \, \stackrel{  }{ \longrightarrow}   \,  {  M   }^{ * }   \, \stackrel{  }{\longrightarrow} \,  {  L   }^{ * }  \,\stackrel{  } {\longrightarrow}  \, 0} {  }

der
\definitionsverweis {Dualräume}{}{}
führt.





}
{} {}

\inputaufgabe
{}
{





Es sei

\mavergleichskette
{\vergleichskette
{a
}
{ \neq }{ 0
}
{  }{
}
{    }{
}
{   }{
}
}
{}{}{}
eine ganze Zahl. Wir betrachten die
\definitionsverweis {kurze exakte Sequenz}{}{}

\mathdisp {0 \longrightarrow \Z \stackrel{\cdot a}{ \longrightarrow} \Z \longrightarrow \Z/( a )  \longrightarrow 0} {  }

von
$\Z$-\definitionsverweis {Moduln}{}{.}
Zeige, dass man die nach

Aufgabe 10.9

exakte Sequenz

\mathdisp {0 \longrightarrow {  ( \Z/( a ) )  }^{ * } \longrightarrow { \Z  }^{ * }  \longrightarrow { \Z  }^{ * }} {  }

bei

\mavergleichskette
{\vergleichskette
{a
}
{ \geq }{ 2
}
{  }{
}
{    }{
}
{   }{
}
}
{}{}{}
nicht nach rechts durch
\mathl{\rightarrow 0}{} exakt fortsetzen kann.





}
{} {}

\inputaufgabe
{}
{





Es sei $R$ ein
\definitionsverweis {kommutativer Ring}{}{}
und sei

\mathdisp {0\stackrel{}{\longrightarrow} L \stackrel{}{\longrightarrow} M
 \mathdisplaybruch\stackrel{}{\longrightarrow} N \stackrel{}{\longrightarrow}0} {  }

eine
\definitionsverweis {kurze exakte Sequenz}{}{}
von
$R$-\definitionsverweis {Moduln}{}{.}
Es gebe ein
$R$-\definitionsverweis {Modul-Erzeugen\-densystem}{}{}
von $L$ mit $k$ Elementen und ein $R$-Modul-Erzeugendensystem von $N$ mit $n$ Elementen. Zeige, dass es ein $R$-Modul-Erzeugendensystem von $M$ mit $k+n$ Elementen gibt.





}
{} {}

\inputaufgabe
{}
{





Es sei $R$ ein
\definitionsverweis {kommutativer Ring}{}{,}
$A$ eine kommutative
\definitionsverweis {endliche}{}{}
$R$-\definitionsverweis {Algebra}{}{}
und $M$ ein
\definitionsverweis {endlicher}{}{}
$A$-\definitionsverweis {Modul}{}{.}
Zeige, dass $M$ auch ein endlicher $R$-Modul ist.





}
{} {}

Die folgenden Aufgaben verwenden den Begriff des artinschen Moduls, der \anfuehrung{dual}{} zum Begriff des noetherschen Moduls ist.





Es sei $R$ ein
\definitionsverweis {kommutativer Ring}{}{.}
Ein
$R$-\definitionsverweis {Modul}{}{}
$M$ heißt \definitionswort {artinsch}{,} wenn jede
absteigende Kette

\mavergleichskettedisp
{\vergleichskette
{ M_1
}
{ \supseteq} { M_2
}
{ \supseteq} { M_3
}
{ \supseteq} { \ldots
}
{ } {
}
}
{}{}{}
von
$R$-\definitionsverweis {Untermoduln}{}{}
\definitionsverweis {stationär}{}{}
wird.






Ein kommutativer Ring $R$ heißt
\definitionswortenp{artinsch}{,} wenn er als $R$-Modul artinsch ist.

\inputaufgabe
{}
{





Es sei $A$ ein
\definitionsverweis {artinscher}{}{}
Integritätsbereich. Man zeige, dass $A$ ein Körper ist. Man gebe ein Beispiel eines artinschen kommutativen Ringes, der kein Körper ist.





}
{} {}

\inputaufgabe
{}
{





Es sei $R$ ein kommutativer Ring und $M$ ein
$R$-\definitionsverweis {Modul}{}{.}
Zeige: Wenn $M$
\definitionsverweis {artinsch}{}{}
und
\maabb {\phi} { M } { M
} {}
$R$-linear und injektiv ist, so ist $\phi$ ein Isomorphismus. Formuliere und beweise auch eine analoge Aussage für den Fall, dass $M$ noethersch ist.





}
{} {}

\inputaufgabegibtloesung
{}
{





Es sei $R$ ein
\definitionsverweis {kommutativer Ring}{}{}
und

\mavergleichskette
{\vergleichskette
{ f
}
{ \in }{ R
}
{  }{
}
{    }{
}
{   }{
}
}
{}{}{}
sei nicht
\definitionsverweis {nilpotent}{}{.}
Zeige, dass es ein
\definitionsverweis {Primideal}{}{}
${\mathfrak p}$ mit

\mavergleichskette
{\vergleichskette
{ f
}
{ \notin }{ {\mathfrak p}
}
{  }{
}
{    }{
}
{   }{
}
}
{}{}{}
gibt.





}
{} {}

\inputaufgabegibtloesung
{}
{





Es sei $\mathfrak a$ ein \definitionsverweis {Radikal}{}{}
in einem kommutativen Ring.
Zeige, dass $\mathfrak a$ der Durchschnitt von
\definitionsverweis {Primidealen}{}{}
ist.





}
{} {Eine Möglichkeit ergibt sich aus der vorstehenden Aufgabe, eine andere aus

Aufgabe 13.5

weiter unten.}

\inputaufgabe
{}
{





Es sei $K$ ein
\definitionsverweis {Körper}{}{} und sei

\mavergleichskette
{\vergleichskette
{ P
}
{ \in }{ K[X]
}
{  }{
}
{    }{
}
{   }{
}
}
{}{}{}
ein nichtkonstantes
\definitionsverweis {Polynom}{}{.}
Zeige, dass $P$ nicht
\definitionsverweis {algebraisch}{}{}
über $K$ ist.





}
{} {}

\inputaufgabe
{}
{





Es sei $K$ ein
\definitionsverweis {Körper}{}{} und

\mavergleichskette
{\vergleichskette
{ L
}
{ = }{ K(X)
}
{  }{
}
{    }{
}
{   }{
}
}
{}{}{}
der
\definitionsverweis {Quotientenkörper}{}{}
des
\definitionsverweis {Polynomrings}{}{}
$K[X]$. Es sei

\mathbed {M} {}
 {K \subseteq M \subseteq L} {}
 {M \neq K} {} {} {,}
ein Zwischenkörper. Zeige, dass

\mavergleichskette
{\vergleichskette
{ M
}
{ \subseteq }{ L
}
{  }{
}
{    }{
}
{   }{
}
}
{}{}{}
eine
\definitionsverweis {endliche Körpererweiterung}{}{}
ist.





}
{} {}

\inputaufgabegibtloesung
{}
{





Es sei $A$ eine
\definitionsverweis {endlich erzeugte}{}{}
$\Z$-\definitionsverweis {Algebra}{}{}
und es sei

\mavergleichskette
{\vergleichskette
{ {\mathfrak m}
}
{ \subseteq }{ A
}
{  }{
}
{    }{
}
{   }{
}
}
{}{}{}
ein
\definitionsverweis {maximales Ideal}{}{.}
Zeige, dass der
\definitionsverweis {Restklassenring}{}{}

\mathl{A/ {\mathfrak m}}{} ein
\definitionsverweis {endlicher Körper}{}{}
ist.





}
{} {}







Es sei $K$ ein
\definitionsverweis {Körper}{}{} und sei $A$ eine kommutative
$K$-\definitionsverweis {Algebra}{}{.}
Man nennt Elemente
\mathl{f_1 , \ldots , f_n \in A}{}
\definitionswort {algebraisch abhängig}{,}
wenn es ein von $0$ verschiedenes Polynom
\mathl{P \in K[X_1 , \ldots , X_n]}{} mit

\mavergleichskettedisp
{\vergleichskette
{P(f_1 , \ldots , f_n)
}
{ =} { 0
}
{ } {
}
{ } {
}
{ } {
}
}
{}{}{}
gibt.







\inputaufgabe
{}
{





Es sei
\mathl{K[X_1 , \ldots , X_n]}{} der
\definitionsverweis {Polynomring}{}{}
über einem
\definitionsverweis {Körper}{}{}
$K$. Zeige, dass die Variablen
\mathl{X_1 , \ldots , X_n}{}
\definitionsverweis {algebraisch unabhängig}{}{}
sind.





}
{} {}

\inputaufgabe
{}
{





Es sei
\mathl{K[X_1 , \ldots , X_n]}{} der
\definitionsverweis {Polynomring}{}{}
über einem
\definitionsverweis {Körper}{}{}
$K$ und seien
\mathl{n+1}{} Polynome

\mavergleichskette
{\vergleichskette
{ f_1 , \ldots , f_{n+1}
}
{ \in }{ K[X_1 , \ldots , X_n]
}
{  }{
}
{    }{
}
{   }{
}
}
{}{}{}
gegeben. Zeige, dass diese
\definitionsverweis {algebraisch abhängig}{}{}
sind.





}
{} {}

\inputaufgabe
{}
{





Es sei
\maabbdisp {\varphi} { { {\mathbb A}_{  K }^{  m   } } } { { {\mathbb A}_{  K }^{  n   } }
} {}
eine
\definitionsverweis {polynomiale Abbildung}{}{}
zwischen
\definitionsverweis {affinen Räumen}{}{}
mit

\mavergleichskette
{\vergleichskette
{m
}
{ < }{ n
}
{  }{
}
{    }{
}
{   }{
}
}
{}{}{.}
Zeige, dass $\varphi$ nicht
\definitionsverweis {surjektiv}{}{}
ist.





}
{} {}

\inputaufgabe
{}
{





Es sei $A$ eine kommutative
$K$-\definitionsverweis {Algebra}{}{}
über einem
\definitionsverweis {Körper}{}{}
$K$ und seien $n$ Elemente

\mavergleichskette
{\vergleichskette
{ f_1 , \ldots ,  f_{n}
}
{ \in }{ A
}
{  }{
}
{    }{
}
{   }{
}
}
{}{}{}
gegeben. Zeige, dass diese Elemente genau dann
\definitionsverweis {algebraisch unabhängig}{}{}
sind, wenn die von diesen Elementen
\definitionsverweis {erzeugte}{}{}
$K$-Algebra
\mathl{K[f_1 , \ldots , f_n ]}{}
\definitionsverweis {isomorph}{}{}
zum
\definitionsverweis {Polynomring}{}{}

\mathl{K[X_1 , \ldots , X_n]}{} ist.





}
{} {}





  \zwischenueberschrift{Aufgaben zum Abgeben}

\inputaufgabe
{3}
{





Es sei $K$ ein
\definitionsverweis {algebraisch abgeschlossener Körper}{}{}
und

\mavergleichskette
{\vergleichskette
{ F
}
{ \in }{ K[X,Y]
}
{  }{
}
{    }{
}
{   }{
}
}
{}{}{}
ein nichtkonstantes Polynom. Zeige, dass der
\definitionsverweis {Restklassenring}{}{}

\mathdisp {K[X,Y]/(F)} {  }

als eine
\definitionsverweis {endliche}{}{}
$K[T]$-Algebra aufgefasst werden kann.





}
{} {}

\inputaufgabe
{3}
{





Es seien
\mathl{R,S,T}{}
\definitionsverweis {kommutative Ringe}{}{}
und seien
\mathl{\varphi:R \rightarrow S}{} und
\mathl{\psi:S \rightarrow T}{} Ringhomomorphismen derart, dass $S$
\definitionsverweis {endlich}{}{}
über $R$ und $T$ endlich über $S$ ist. Zeige, dass dann auch $T$ endlich über $R$ ist.





}
{} {}

\inputaufgabe
{5}
{





Es sei $A$ ein
\definitionsverweis {kommutativer Ring}{}{}
und sei

\mathdisp {0 \longrightarrow M \longrightarrow N \longrightarrow P \longrightarrow 0} {  }

eine
\definitionsverweis {kurze exakte Sequenz}{}{}
von $A$-Moduln. Man zeige, dass $N$ genau dann
\definitionsverweis {artinsch}{}{}
ist, wenn $M$ und $P$ artinsch sind.





}
{} {}

\inputaufgabe
{4 (1+3)}
{





Es sei $R$ ein
\definitionsverweis {kommutativer Ring}{}{}
und $M_i$,

\mavergleichskette
{\vergleichskette
{ i
}
{ \in }{ \N
}
{  }{
}
{    }{
}
{   }{
}
}
{}{}{,}
seien $R$-Moduln mit fixierten
$R$-\definitionsverweis {Modulhomomorphismen}{}{}
\maabbdisp {\varphi_i} { M_i } { M_{i+1}
} {.}
Die Sequenz

\mathdisp {\ldots \longrightarrow M_i \longrightarrow M_{i+1} \longrightarrow M_{i+2} \longrightarrow M_{i+3} \longrightarrow \ldots} {  }

heißt
\definitionswortenp{exakt}{,} wenn für alle $i$ gilt, dass

\mavergleichskette
{\vergleichskette
{ \operatorname{Kern} \,(\varphi_{i})
}
{ = }{ \operatorname{Bild} \, (\varphi_{i-1})
}
{  }{
}
{    }{
}
{   }{
}
}
{}{}{}
ist.
\aufzaehlungzwei {Zeige, dass diese Definition im Falle einer kurzen exakten Sequenz mit der

Definition 10.2
 in der Vorlesung
übereinstimmt.
} {Es sei nun

\mavergleichskette
{\vergleichskette
{ R
}
{ = }{ K
}
{  }{
}
{    }{
}
{   }{
}
}
{}{}{}
ein Körper, die $M_i$ seien  endlich erzeugt,

\mavergleichskette
{\vergleichskette
{ M_0
}
{ = }{ 0
}
{  }{
}
{    }{
}
{   }{
}
}
{}{}{}
und alle

\mavergleichskette
{\vergleichskette
{ M_i
}
{ = }{ 0
}
{  }{
}
{    }{
}
{   }{
}
}
{}{}{}
für

\mavergleichskette
{\vergleichskette
{ i
}
{ \geq }{ n
}
{  }{
}
{    }{
}
{   }{
}
}
{}{}{}
für ein gewisses $n$. Zeige, dass

\mavergleichskettedisp
{\vergleichskette
{ \sum_{i = 0}^{n} (-1)^i \operatorname{dim}_K M_i
}
{ =} { 0
}
{ } {
}
{ } {
}
{ } {
}
}
{}{}{.}
}





}
{} {}

\inputaufgabe
{3}
{





Es sei $K$ ein
\definitionsverweis {Körper}{}{} und $A$ eine
\definitionsverweis {endliche}{}{}
$K$-Algebra. Zeige: Dann ist $A$
\definitionsverweis {artinsch}{}{.}





}
{} {}
